A bare address prints its underscores and tildes: \href{https://example.com/a_b_c}{https://example.com/a\_b\_c} and
\href{https://example.com/~me/page}{https://example.com/\textasciitilde{}me/page}, \href{https://example.com/x_y~z}{https://example.com/x\_y\textasciitilde{}z} in angle brackets,
and \href{http://www.example.com/w_x}{www.example.com/w\_x} with no scheme.

A link written out keeps its URL's percent sign inside a footnote,\footnote{See \href{https://example.com/n_o~p}{https://example.com/n\_o\textasciitilde{}p} and \href{https://example.com/q\%20r}{the page}.} in bold,
\textbf{\href{https://example.com/u\%20v}{the bold link}}, in bold italics,
\textbf{\emph{\href{https://example.com/w\%20x}{the other}}}, and in a heading:

\subsection{Heading \href{https://example.com/h\%20i}{linked} and \href{https://example.com/j_k}{https://example.com/j\_k}}

\begin{itemize}
\item \href{https://example.com/l_m~n}{https://example.com/l\_m\textasciitilde{}n} in a list
\end{itemize}

A written-out link's underscore and tilde are as they were:
\href{https://example.com/a_b~c}{text}.

